% !PRINTED 701
\weberpage{701}
%<<CONT
and therefore
\weq{20}{\mathop{\mathrm{Lim}}\limits_{s=1}\left(S - \frac{\gamma_s}{s-1}\right) = -\gamma\,\Gamma'(1) + D_1 = C_1,}
% [sic] The source prints gamma_s, +D_1, and =C_1. With the intended gamma^s, (19) gives the limit -gamma*Gamma'(1)-D_1. Definition (17), which subtracts gamma/(s-1), instead gives C_1=-gamma*Gamma'(1)+gamma*log(gamma)-D_1. The printed subscript, sign, and identification are retained and distinguished here.
is finite, as was to be proved\footnote{The theorems on $\Gamma$-functions used here are to be found in the more detailed textbooks of the integral calculus, e.\,g. \gesperrt{Serret-Harnack}, Lehrbuch der Differential- und Integralrechnung, Bd.\ II, S.\,170\,f. (Leipzig 1885).}.

\begin{center}
\S.\,192.\\[4pt]
\gesperrt{Application to the determination of the class number.}
\end{center}

We now make application of these theorems to the theory of the field $\Omega$. By the numbers $\mu_n$ of Theorem 5. we understand the norms of all the ideals of the field $\Omega$, and hence by $T$ the number of ideals in $\Omega$ whose norm is not greater than a given positive quantity $t$, and we obtain
\weq{1}{\mathop{\mathrm{Lim}}\limits_{s=1}\sum \frac{s-1}{N(\mathfrak{a})^s} = \mathop{\mathrm{Lim}}\limits_{t=\infty}\frac{T}{t},}
and in this the sum on the left side extends over \gesperrt{all} ideals $\mathfrak{a}$ of the field. The ideals now split, according to \S.\,153, into a finite number of classes
\[
A_1, A_2, \ldots, A_h,
\]
and the great goal is the determination of this number $h$, the \gesperrt{class number}.

The ideals $\mathfrak{a}_1$ of one of these classes $A_1$ are characterised by this, that their products with one and the same integral functional $\varphi_1$, which belongs to the class $A_1^{-1}$, are \gesperrt{principal ideals}, so that therefore
\[
\varphi_1 \mathfrak{a}_1 = \alpha
\]
is an integer of the field $\Omega$.

If the norm of $\mathfrak{a}_1$ is not greater than $t$, then
\[
N_a(\alpha) \leqq N_a(\varphi_1)\, t = t_1.
\]
If $T_1$ is the number of ideals of the class $A_1$ whose norms are not greater than $t$, then $T_1$ is at the same time the number of the non-
%CONT>>
